\documentclass[a4paper]{article}

% Basic packages
\usepackage[fontset=fandol]{ctex}
\usepackage{amsmath, amssymb}
\usepackage{graphicx}
\usepackage[margin=2.5cm]{geometry}
\usepackage[most]{tcolorbox}
\usepackage{etoolbox}
\usepackage{listings}
\usepackage{hyperref}
\usepackage{booktabs}
\usepackage{subcaption}
\usepackage{float}
\usepackage{tikz}

% ===== Highlight boxes =====
\newtcolorbox{knowledgebox}[1]{
    enhanced, colback=blue!5!white, colframe=blue!75!black, colbacktitle=blue!75!black,
    coltitle=white, fonttitle=\bfseries, title=#1, attach boxed title to top left={yshift=-2mm, xshift=2mm},
    boxrule=1pt, sharp corners
}

\newtcolorbox{importantbox}[1]{
    enhanced, colback=yellow!10!white, colframe=yellow!80!black, colbacktitle=yellow!80!black,
    coltitle=black, fonttitle=\bfseries, title=#1, sharp corners
}

\newtcolorbox{warningbox}[1]{
    enhanced, colback=red!5!white, colframe=red!75!black, colbacktitle=red!75!black,
    coltitle=white, fonttitle=\bfseries, title=#1, sharp corners
}

% ===== Code listings =====
\lstset{
    basicstyle=\ttfamily\small,
    keywordstyle=\color{blue},
    stringstyle=\color{red!60!black},
    commentstyle=\color{green!60!black},
    breaklines=true,
    frame=single,
    numbers=left,
    numberstyle=\tiny\color{gray},
    captionpos=b,
    extendedchars=false
}

% ===== Front-page metadata =====
\newcommand{\notetitle}{机器学习任务攻略}
\newcommand{\notesubtitle}{李宏毅《机器学习》2025 · 第二节（General Guidance）}
\newcommand{\noteauthors}{基于李宏毅老师课程整理}
\newcommand{\notedate}{\today}
\newcommand{\videochannel}{李宏毅深度学习课堂（Bilibili 搬运）}
\newcommand{\videoduration}{约 51 分钟}
\newcommand{\videourl}{https://www.bilibili.com/video/BV1TAtwzTE1S?p=3}

\begin{document}

\pagestyle{empty}

% ===== Cover page =====
\begin{center}
    \vspace*{1.5cm}
    {\Huge\bfseries \notetitle\par}
    \vspace{0.4cm}
    {\LARGE \notesubtitle\par}
    \vspace{0.8cm}
    \includegraphics[width=0.62\textwidth]{video.jpg}\par
    \vspace{0.8cm}
    {\large \noteauthors\par}
    {\large \notedate\par}
    \vspace{0.8cm}
    \begin{tcolorbox}[width=0.8\textwidth,center,sharp corners,
                      colback=black!3!white,colframe=black!50!white]
        \centering
        \textbf{视频信息}\\[3pt]
        频道：\videochannel \\
        时长：\videoduration \\
        链接：\href{\videourl}{\videourl}
    \end{tcolorbox}
\end{center}

\newpage

% ===== TOC =====
\tableofcontents
\newpage

\pagestyle{plain}

% =====================================================================
\section{引言：作业通关的总攻略}
% =====================================================================

从这一节开始，李宏毅老师不再讲单一模型，而是给出一套\textbf{适用于前期所有作业的通用攻略}（General Guidance）。老师把它比喻成游戏里的开局神装——``就跟魔关羽一样，开局就送，可以帮你打赢前期所有的副本''。

\subsection{所有作业的共同格式}

之后好几个作业看起来都大同小异：你会拿到一堆\textbf{训练资料}，每笔含输入 $x$ 与对应的标签 $\hat{y}$；而\textbf{测试资料}只有 $x$、没有 $\hat{y}$，需要你的模型去预测。

\begin{knowledgebox}{各作业的 $X \to Y$ 一览}
    \begin{itemize}
        \item \textbf{作业 1（回归）}：$X=$ 过去的疫情/观看数据，$Y=$ 未来某天的数值
        \item \textbf{作业 2（语音辨识）}：$X=$ 一小段声音讯号，$Y=$ 对应的音素（phoneme，可暂时理解为 KK 音标）
        \item \textbf{作业 3（影像辨识）}：$X=$ 一张图片，$Y=$ 图片里是什么东西
        \item \textbf{作业 4（语者辨识）}：$X=$ 一段声音讯号，$Y=$ 说话的是哪一个人（银行客服的身份验证就用这种系统）
        \item \textbf{作业 5（机器翻译）}：$X=$ 一种语言的句子，$Y=$ 另一种语言的句子
    \end{itemize}
\end{knowledgebox}

\subsection{训练的三个步骤（回顾）}

训练资料拿来训练 model，过程就是上一节讲的三个步骤：

\begin{figure}[H]
    \centering
    \includegraphics[width=0.82\textwidth]{figures/fig02_framework.jpg}
    \caption{机器学习框架：Step 1 写出带未知参数的函数 $y=f_{\boldsymbol{\theta}}(x)$；Step 2 从训练资料定义损失 $L(\boldsymbol{\theta})$；Step 3 解最佳化问题求 $\boldsymbol{\theta}^*=\arg\min_{\boldsymbol{\theta}} L$。最后用 $f_{\boldsymbol{\theta}^*}$ 去标注测试资料。}
    \label{fig:framework}
    \footnotetext{视频画面时间区间：00:03:50--00:04:10。}
\end{figure}

\begin{enumerate}
    \item 写出一个有未知数的函数 $y=f_{\boldsymbol{\theta}}(x)$，其中 $\boldsymbol{\theta}$ 代表 model 里所有未知参数，输入 $x$ 称为 feature
    \item 定义损失函数 $L(\boldsymbol{\theta})$，衡量一组参数有多好/多坏
    \item 解最佳化问题，找出让损失最小的 $\boldsymbol{\theta}^*$，再把它代回函数用于测试资料
\end{enumerate}

直接执行助教的 sample code，往往只能通过 simple baseline。\textbf{想做得更好，就要按下面这张攻略图逐步诊断。}

\subsection{攻略全景图：一棵诊断决策树}

整套攻略的核心，是一棵\textbf{从损失出发的决策树}。后续每一节都是在走这棵树的某一条分支，建议先把它印在脑海里：

\begin{figure}[H]
    \centering
    \includegraphics[width=0.95\textwidth]{figures/fig01_general_guide.jpg}
    \caption{机器学习任务攻略决策树（General Guide）：先看\textbf{训练资料}上的损失，再看\textbf{测试资料}上的损失，沿不同分支对症下药。}
    \label{fig:guide}
    \footnotetext{视频画面时间区间：00:48:45--00:49:15（此为讲完全部情形后的完整版攻略图）。}
\end{figure}

\begin{importantbox}{攻略的核心顺序：先看训练损失，再看测试损失}
    很多同学一看到 Kaggle 上（测试资料）结果不好，就直接下结论说``overfitting''。\textbf{这是最常见的错误。} 正确的第一步永远是：

    \begin{center}
    \fbox{先检查 \textbf{训练资料} 的 loss，确认 model 在训练资料上``学起来''了，再去看测试资料。}
    \end{center}

    决策树主干：
    \begin{itemize}
        \item \textbf{训练 loss 大} $\rightarrow$ 是 Model Bias 还是 Optimization 失败？（第 2 节）
        \item \textbf{训练 loss 小、测试 loss 大} $\rightarrow$ 是 Overfitting 还是 Mismatch？（第 3、5 节）
        \item \textbf{训练 loss 小、测试 loss 也小} $\rightarrow$ 恭喜通关（\texttt{:-)}）
    \end{itemize}
\end{importantbox}

% =====================================================================
\section{训练损失大：Model Bias 还是 Optimization 失败}
% =====================================================================

如果训练资料上的 loss 就很大，说明 model 连训练资料都没学好。原因有两个：\textbf{Model Bias} 与 \textbf{Optimization 失败}。

\subsection{Model Bias：模型太简单，针不在海里}

\begin{figure}[H]
    \centering
    \includegraphics[width=0.92\textwidth]{figures/fig03_model_bias.jpg}
    \caption{Model Bias：model 能表示的所有函数构成一个集合（function set），但这个集合``太小''，里面根本没有能让 loss 变低的函数（橙点不在蓝色集合内）。解法是重新设计更有弹性的 model。}
    \label{fig:model_bias}
    \footnotetext{视频画面时间区间：00:06:00--00:07:45。}
\end{figure}

把不同参数 $\boldsymbol{\theta}$ 代入函数，会得到不同的函数；把它们全部集合起来，就是这个 model 能表示的 function set。\textbf{Model Bias 指的是：这个集合太小，里面没有任何一个函数能让 loss 足够低。}

\begin{knowledgebox}{大海捞针的比喻}
    寻找低 loss 的函数，就像\textbf{在大海里捞针}。Model Bias 的情况是：\textbf{针根本不在海里}。所以无论 optimization 多努力，找到的 $\boldsymbol{\theta}^*$ 也只是这些差函数里``最好的那一个''——``它只是鲁蛇里面的霸主，还是一个鲁蛇''，loss 依然不够低。
\end{knowledgebox}

\textbf{解法：重新设计、给 model 更大的弹性。} 承接上一节的做法：

\[
y = b + w x_1
\;\xrightarrow{\ \text{增加 feature}\ }\;
y = b + \sum_{j=1}^{56} w_j x_j
\;\xrightarrow{\ \text{Deep Learning}\ }\;
y = b + \sum_i c_i\, \sigma\!\Big(b_i + \sum_j w_{ij} x_j\Big)
\]

即：增加输入的 feature 数量、或用更多神经元/更多层的深度网络来扩大 function set。

\begin{warningbox}{再次提醒：Model Bias 的 ``bias'' 不是参数 $b$}
    Model Bias 是``模型结构的局限''，与函数里那个截距参数 $b$（bias）完全是两回事。这一点上一节已强调过，这里再次确认，避免混淆。
\end{warningbox}

\subsection{Optimization 失败：针在海里，却捞不出来}

并不是训练 loss 大就一定是 Model Bias。另一种可能是 \textbf{optimization 没做好}。本课程的 optimization 方法只用 gradient descent，而它有很多问题——比如会卡在 local minima，找不到真正能让 loss 很低的参数。

\begin{knowledgebox}{两种情况的比喻对比}
    \begin{itemize}
        \item \textbf{Model Bias}：针不在海里（function set 里没有好函数）$\rightarrow$ 白忙一场
        \item \textbf{Optimization 失败}：针在海里，但 gradient descent 没本事把它捞出来（好函数存在，却找不到）
    \end{itemize}
\end{knowledgebox}

\subsection{如何区分二者：比较不同的模型}

看到训练 loss 大，怎么知道是``海里没针''还是``捞不出针''？\textbf{建议的方法：比较不同复杂度的模型。}

\subsubsection{ResNet 的经典实验}

\begin{figure}[H]
    \centering
    \includegraphics[width=0.82\textwidth]{figures/fig04_resnet.jpg}
    \caption{ResNet 论文（2015，arXiv:1512.03385）的实验：在\textbf{测试资料}上，56 层网络（红）的 error 反而比 20 层（黄）更高。}
    \label{fig:resnet}
    \footnotetext{视频画面时间区间：00:10:30--00:11:30。}
\end{figure}

当年很多人看到这张图就说``56 层太深、overfitting 了，所以深度学习不 work''。\textbf{但这并不是 overfitting。} 检查\textbf{训练资料}会发现：56 层在训练资料上的 loss 也比 20 层高。

\begin{importantbox}{为什么这是 Optimization 失败，而非 Model Bias 或 Overfitting}
    \begin{itemize}
        \item \textbf{不是 Overfitting}：因为它连\textbf{训练}资料的 loss 都更高（overfitting 是训练好、测试差）
        \item \textbf{不是 Model Bias}：56 层的弹性\textbf{一定} $\geq$ 20 层。56 层只要让前 20 层和那个 20 层网络一样、后 36 层全做 identity（原样输出），就能复制 20 层的结果
        \item \textbf{所以是 Optimization 失败}：56 层有能力做得更好，gradient descent 却没找到那组参数
    \end{itemize}
\end{importantbox}

\subsubsection{实用建议与观看人数例子}

\begin{knowledgebox}{诊断 Optimization 的实用建议}
    遇到一个没做过的新问题，可以\textbf{先跑比较浅、比较小的网络，甚至非深度学习的方法}（如 linear model、SVM）。这些模型``会竭尽全力在能力范围内找出最好的参数''，比较不会有 optimization 失败的问题，可以作为 baseline。

    之后再跑深的 model：如果\textbf{深的 model 弹性更大、loss 却压不过浅的 model}，那就说明 optimization 出了问题。
\end{knowledgebox}

上一节的观看人数预测实验正是一例（训练资料为 2017--2020 年）：

\begin{table}[H]
    \centering
    \caption{不同层数网络在训练资料上的 loss（单位：千人次）}
    \label{tab:layers}
    \begin{tabular}{lccccc}
        \toprule
        \textbf{层数} & 1 层 & 2 层 & 3 层 & 4 层 & \textbf{5 层} \\
        \midrule
        \textbf{训练 loss} & 0.28k & 0.18k & 0.14k & 0.10k & \textbf{0.34k} \\
        \bottomrule
    \end{tabular}
\end{table}

1$\to$4 层 loss 一路递减，到 \textbf{5 层却暴增到 0.34k}。这显然不是 Model Bias（4 层都能做到 0.10k，弹性更大的 5 层没理由更差），而是 \textbf{optimization 没做好}。至于 optimization 失败具体怎么解决，留待\textbf{下一节课}（local minima 与 saddle point）。

% =====================================================================
\section{训练损失小、测试损失大：Overfitting}
% =====================================================================

假设你已经让训练 loss 变小了，接着看测试 loss。若测试 loss 也小、且小于 strong baseline，那就通关了。\textbf{但若训练 loss 小、测试 loss 大，这才可能是真正的 Overfitting。}

\begin{warningbox}{Overfitting 的严格定义}
    \textbf{训练 loss 小} $+$ \textbf{测试 loss 大} $=$ Overfitting。两个条件缺一不可。老师说：``十个同学有八个''被问到``训练 loss 多少''时都答不上来——请务必先记录训练 loss，确认 optimization 没问题、model 够大之后，再来谈 testing 的问题。
\end{warningbox}

\subsection{为什么会发生 Overfitting}

\begin{figure}[H]
    \centering
    \includegraphics[width=0.85\textwidth]{figures/fig05_overfitting.jpg}
    \caption{Overfitting 直觉：真实关系是虚线（不可观测），训练资料只是其上 sample 出的几个蓝点。弹性极大的 model 会精确穿过训练点，但在没有数据约束处``freestyle''乱长，导致测试点（橙）上误差很大。}
    \label{fig:overfitting}
    \footnotetext{视频画面时间区间：00:19:00--00:20:15。}
\end{figure}

\begin{knowledgebox}{一个极端的``一无是处''函数}
    设想一个很废的函数：输入 $x$ 时，先查它有没有出现在训练资料里——有就输出对应的 $\hat{y}$，没有就输出一个\textbf{随机值}。

    这个函数啥也没学到，但它在\textbf{训练资料上 loss $=0$}（全背下来了），在\textbf{测试资料上 loss 极大}（全靠乱猜）。这就是 overfitting 的极端形态。
\end{knowledgebox}

一般情况下，当 model 弹性（自由度）很大时，只用少数几个训练点作约束，其余地方它可以产生各式各样奇怪的曲线，于是训练好、测试差。

\subsection{解法一：增加训练资料 / Data Augmentation}

\begin{figure}[H]
    \centering
    \includegraphics[width=0.85\textwidth]{figures/fig06_data_aug.jpg}
    \caption{增加训练资料：蓝点变多后，即使 model 弹性很大，也会被密集的数据``限制住''，形状被迫贴近真实的二次曲线。}
    \label{fig:data_aug}
    \footnotetext{视频画面时间区间：00:21:00--00:22:45。}
\end{figure}

\textbf{增加训练资料往往是最有效的解法。} 但作业里\textbf{不允许自己额外收集资料}（以免大家把精力花在搜集资料而非机器学习核心技术上）。作业里能用的是 \textbf{Data Augmentation}：

\begin{importantbox}{Data Augmentation（资料增强）}
    用你\textbf{对问题的理解}，从已有资料创造出新资料——这不算使用额外资料。例如影像辨识中把图片\textbf{左右翻转}、或\textbf{截出一块放大}，资料量就翻倍。

    \textbf{但要 augment 得有道理：} 影像辨识中很少把图片\textbf{上下颠倒}，因为那不是真实世界会出现的画面，机器看了反而会学到奇怪的东西。
\end{importantbox}

\subsection{解法二：限制模型（但别限制过度）}

另一个方向是\textbf{不要给 model 那么大的弹性、给它一些限制}，最好让 model 正好契合产生资料背后的规律（例如已知是二次曲线，就限制 model 只能是二次曲线）。

\begin{knowledgebox}{给 model 制造限制的常见手段}
    \begin{itemize}
        \item 给\textbf{比较少的参数} / 比较少的神经元
        \item \textbf{共用参数}（让一些参数取相同数值）——CNN 就是典型
        \item 用\textbf{比较少的 feature}
        \item \textbf{Early Stopping}、\textbf{Regularization}、\textbf{Dropout}（后续课程详讲；作业一的 sample code 里已含 early stopping）
    \end{itemize}
\end{knowledgebox}

\begin{knowledgebox}{为什么 CNN 在影像上反而更好？}
    Fully-connected network 是\textbf{弹性较大}的架构，它能表示的 function set 较大；CNN 是\textbf{受限}的架构，function set 较小、且包含于 fully-connected 之内。但正因为 CNN 针对影像特性施加了\textbf{合理的限制}，它在影像问题上反而做得更好。``弹性大''不等于``更好''。
\end{knowledgebox}

\begin{figure}[H]
    \centering
    \includegraphics[width=0.85\textwidth]{figures/fig07_constrain.jpg}
    \caption{限制过度的反例：若强行规定 model 只能是直线 $y=a+bx$，则没有任何直线能同时通过这三个点，测试结果同样很差——但这不是 overfitting，而是\textbf{回到了 Model Bias}。}
    \label{fig:constrain}
    \footnotetext{视频画面时间区间：00:28:00--00:29:00。}
\end{figure}

\begin{warningbox}{限制不能给太多——否则回到 Model Bias}
    给的限制太多、model 太死板，就会重新陷入 Model Bias：结果不好不是因为 overfitting，而是因为 function set 又变得太小、装不下好函数。
\end{warningbox}

\subsection{Bias-Complexity Trade-off}

于是出现一个矛盾：model 太简单有 Model Bias，太复杂有 Overfitting。

\begin{figure}[H]
    \centering
    \includegraphics[width=0.78\textwidth]{figures/fig08_tradeoff.jpg}
    \caption{Bias-Complexity Trade-off：随 model 复杂度（更多 feature/参数）增加，训练 loss 单调下降；而测试 loss 先降后升，呈 U 形。最优点在中间``不偏不倚''处。}
    \label{fig:tradeoff}
    \footnotetext{视频画面时间区间：00:29:40--00:30:30。}
\end{figure}

\begin{importantbox}{复杂度与损失的关系}
    \begin{itemize}
        \item \textbf{训练 loss}：随复杂度增加\textbf{单调下降}
        \item \textbf{测试 loss}：先随复杂度下降，超过某临界点后因 overfitting \textbf{突然回升}（U 形）
    \end{itemize}
    我们想选一个\textbf{中庸的 model}——不太复杂也不太简单，恰好在测试 loss 最低处。怎么选？见下一节。

    （``复杂度/弹性''在本节只是直观描述，其严格定义将在后续``吴沛元老师''的课程中给出。）
\end{importantbox}

% =====================================================================
\section{模型选择（Model Selection）}
% =====================================================================

\subsection{为什么不能用 public 测试集来挑模型}

最直觉的做法：把几个模型的结果都传上 Kaggle，看哪个分数最低就选哪个。\textbf{老师强烈不建议这么做。}

\begin{warningbox}{一兆个废模型的极端例子}
    假设有一兆个``一无是处''的模型，每个在测试资料上都是\textbf{随机输出}。把它们的结果全传上 Kaggle，总会有某一个（比如编号 56789）\textbf{碰巧}在 public 测试集上得到好分数。你以为它是好 model，其实它只是随机蒙中。

    后果：\textbf{public 排前十，private 一公布就崩}——掉到 300 名甚至 1000 名之外。``这件事每年都会发生，并不是传说。''
\end{warningbox}

\begin{knowledgebox}{Benchmark 上``超越人类''的真相}
    这也解释了为什么在公开的 benchmark corpus（如语音辨识的 LibriSpeech、Switchboard）上，机器常``超越人类''——因为测试集是\textbf{公开的}，只要刷得够多，再废的模型也能在上面蒙到好成绩。2016 年微软、IBM 宣称语音辨识超越人类，就是 benchmark 上的结果；现实生活中（private 场景）显然没有。这类宣传``比较像是骗骗麻瓜的商业实力''。
\end{knowledgebox}

为此，Kaggle 把测试集分成 public 与 private，且\textbf{限制每日上传次数}，就是为了防止大家靠不断随机上传来碰运气。

\subsection{正确做法：切出 Validation Set}

\begin{figure}[H]
    \centering
    \includegraphics[width=0.85\textwidth]{figures/fig09_validation.jpg}
    \caption{把训练资料再切成 Training Set 与 Validation Set（如 90\% / 10\%）。在 Training Set 上训练、在 Validation Set 上衡量并\textbf{挑选}模型，再上传。}
    \label{fig:validation}
    \footnotetext{视频画面时间区间：00:40:45--00:41:30。}
\end{figure}

\begin{importantbox}{用 Validation Set 挑模型}
    用 validation set 上的分数来选 model，则 public 测试集的分数就能较好地反映 private 测试集的分数，不会出现``public 极好、private 极差''的悲剧。

    \textbf{理想做法：} 用 validation loss 最小的那个直接定案，\textbf{尽量不要再根据 public 分数去反复调整模型}——看 public 结果调得越多，就越可能 fit 在 public 上，private 又会崩。
\end{importantbox}

\subsection{N-fold 交叉验证}

如果担心``万一 validation set 恰好分得很差''，可以用 \textbf{N-fold Cross Validation}。

\begin{figure}[H]
    \centering
    \includegraphics[width=0.90\textwidth]{figures/fig10_nfold.jpg}
    \caption{N-fold 交叉验证（此处 $N=3$）：训练资料切成 3 份，轮流让其中一份当 Validation、其余当 Training，共 3 种切法。每个模型在 3 种切法下各跑一次，取平均 mse；平均最低者（此处 Model 1，avg mse $=0.3$）即为最佳，再用\textbf{全部}训练资料重训后用于测试集。}
    \label{fig:nfold}
    \footnotetext{视频画面时间区间：00:45:10--00:46:30。}
\end{figure}

\begin{enumerate}
    \item 把训练资料切成 $N$ 等份（例 3 份）
    \item 轮流取一份作 validation、其余作 training，得到 $N$ 种 train/val 配置
    \item 每个候选模型在这 $N$ 种配置下都跑一遍，把结果\textbf{平均}
    \item 选平均结果最好的模型，再用\textbf{全部训练资料}重新训练，最后用于测试集
\end{enumerate}

% =====================================================================
\section{回到观看人数预测：Mismatch}
% =====================================================================

\subsection{2 月 26 日的预测结果}

上一节大家投票多数选了 3 层网络，就直接用它预测 2021 年的观看人数：

\begin{figure}[H]
    \centering
    \includegraphics[width=0.82\textwidth]{figures/fig11_predict226.jpg}
    \caption{3 层网络预测结果：红线为真实值、蓝线为预测值。模型在 2/26（2021 年观看人数最高的一天）预测得非常惨，误差高达 $e=2.58$k。}
    \label{fig:predict}
    \footnotetext{视频画面时间区间：00:46:54--00:48:15。}
\end{figure}

老师还跑了 1、2、4 层对比：2 层、3 层误差都是 2 点多 k，1 层和 4 层反而好些（约 1.8k）。但\textbf{四个模型不约而同都以为 2/26 是个低点}——因为根据历史，礼拜五观看人数通常最少（``礼拜五晚上大家出去玩了''）。可 2/26 偏偏反常地成了峰值。

\begin{knowledgebox}{这不能怪模型}
    模型只看得到过去的观看数据，并不知道``那天有什么特别''。这种错误既不是 model bias，也不是单纯的 overfitting，而是另一种形式——\textbf{Mismatch}。
\end{knowledgebox}

\subsection{Mismatch：训练与测试分布不同}

\begin{figure}[H]
    \centering
    \includegraphics[width=0.85\textwidth]{figures/fig12_mismatch.jpg}
    \caption{Mismatch 范例（作业 11）：训练资料是真实照片，测试资料却是涂鸦线条画，两者分布完全不同。此时单纯增加训练资料\textbf{也无济于事}。}
    \label{fig:mismatch}
    \footnotetext{视频画面时间区间：00:49:00--00:51:00。}
\end{figure}

\begin{importantbox}{Mismatch vs Overfitting}
    \begin{itemize}
        \item \textbf{Overfitting}：可以靠\textbf{增加训练资料}来缓解
        \item \textbf{Mismatch}：训练资料与测试资料的\textbf{分布本身不一样}。此时再增加训练资料也没用——因为多出来的还是``分布不对''的资料
    \end{itemize}
\end{importantbox}

大多数作业都已设计好、不会遇到 mismatch，\textbf{唯一的例外是作业 11}（专为 mismatch 而设计：训练真实照片、测试涂鸦）。

\begin{knowledgebox}{COVID 数据的 mismatch 隐患}
    以作业一为例：若用 2020 年资料当训练、2021 年资料当测试，由于两年背后分布差很多，``怎么训练、什么模型都会铲掉''。所以助教后来改用了别的方式来切分训练/测试资料。
\end{knowledgebox}

\textbf{怎么判断是不是 mismatch？} 没有公式可套，要靠你\textbf{对资料本身、对训练/测试资料产生方式的理解}来判断。

% =====================================================================
\section{总结与延伸}
% =====================================================================

本节给出了一套贯穿前期所有作业的\textbf{诊断决策树}：从损失出发，逐层定位问题、对症下药。

\begin{importantbox}{攻略决策树（文字版复盘）}
    \textbf{第一步——看训练资料的 loss：}
    \begin{itemize}
        \item \textbf{大}：
            \begin{itemize}
                \item Model Bias $\rightarrow$ 把 model 变大（加 feature、加深）
                \item Optimization 失败 $\rightarrow$ 换更好的最佳化方法（下一节）
            \end{itemize}
        \item \textbf{小} $\rightarrow$ 第二步——看测试资料的 loss：
            \begin{itemize}
                \item 小 $\rightarrow$ 通关
                \item 大：Overfitting $\rightarrow$ 增数据 / Data Augmentation / 限制模型；或 Mismatch $\rightarrow$ 分布问题，增数据无用
            \end{itemize}
    \end{itemize}
\end{importantbox}

\begin{knowledgebox}{四种``结果不好''的根因对照}
    \begin{center}
    \begin{tabular}{lccl}
        \toprule
        \textbf{问题} & \textbf{训练 loss} & \textbf{测试 loss} & \textbf{对策} \\
        \midrule
        Model Bias    & 大 & 大 & 加大模型（feature/深度） \\
        Optimization 失败 & 大 & 大 & 更好的最佳化（下一节） \\
        Overfitting   & 小 & 大 & 增数据/增强、限制模型 \\
        Mismatch      & 小 & 大 & 对齐分布（增数据无效） \\
        \bottomrule
    \end{tabular}
    \end{center}
    注：Model Bias 与 Optimization 失败都表现为``训练 loss 大''，靠\textbf{比较不同模型}来区分；Overfitting 与 Mismatch 都表现为``训练好测试差''，靠\textbf{对资料分布的理解}来区分。
\end{knowledgebox}

\begin{knowledgebox}{本节关键概念清单}
    \begin{center}
    \begin{tabular}{lll}
        \toprule
        \textbf{概念} & \textbf{英文} & \textbf{说明} \\
        \midrule
        模型偏差 & Model Bias & function set 太小，装不下好函数 \\
        最佳化失败 & Optimization Issue & 好函数存在，gradient descent 找不到 \\
        过拟合 & Overfitting & 训练好、测试差（弹性过大） \\
        资料增强 & Data Augmentation & 用领域理解从已有资料造新资料 \\
        偏差-复杂度权衡 & Bias-Complexity Trade-off & 测试 loss 随复杂度呈 U 形 \\
        验证集 & Validation Set & 从训练资料切出，用于挑模型 \\
        交叉验证 & N-fold Cross Validation & 轮流当验证集，平均后挑模型 \\
        公开/私有测试集 & Public / Private Testing Set & 防止靠刷 public 碰运气 \\
        不匹配 & Mismatch & 训练与测试分布不同 \\
        \bottomrule
    \end{tabular}
    \end{center}
\end{knowledgebox}

\subsection{下一讲预告}

\begin{itemize}
    \item \textbf{Optimization 为什么会失败、怎么办}：局部最小值（local minima）与鞍点（saddle point）
    \item 之后还会讲到 batch 与 momentum、自动调整学习率、损失函数的选择等``类神经网络训练不起来怎么办''系列
    \item Regularization、Dropout、CNN 等``限制模型''手段的细节，也会在后续课程展开
\end{itemize}

\vspace{1cm}

\begin{center}
    \rule{0.5\textwidth}{0.5pt}
    \vspace{0.3cm}

    \textbf{本节完 · 下一节：类神经网络训练不起来怎么办（一）—— Local Minima 与 Saddle Point}

    课程视频：\href{\videourl}{\videourl}

    \vspace{0.3cm}
    \rule{0.5\textwidth}{0.5pt}
\end{center}

\end{document}
